% P01: replaces the FIRST paragraph of
% \subsection{Model Conversion and Actuator Coupling}\label{sec:sns_existing_tools}
% EXISTING STARTS: 'MyoConverter supplies the initial OpenSim-to-MuJoCo translation. I added procedures...'
% EXISTING ENDS:   '...exchange signals at a \qty{2}{\milli\second} step.'
MyoConverter supplies the initial OpenSim-to-MuJoCo translation, and the conversion checks that accompanied it found and corrected five concrete problems. My first attempt converted the robot model itself, in which the braided-pneumatic routing had been laid over the human Gait2392 skeleton with the human muscle parameters left unchanged, and MyoConverter's equilibrium checks flagged the combination as infeasible because the routed path lengths depart far from the lengths those parameters describe. The workable sequence converted the unmodified human reference model first, then regenerated the robot routing on both legs through a mirroring procedure that also restored muscles disabled by earlier test configurations. Three very small hip rotators per side, the quadratus femoris, gemelli, and piriformis, are poorly conditioned even in OpenSim, and I rendered their actuators inert at model load, so the converted model keeps its 92 muscle-tendon actuators with six of them carrying no force. MuJoCo's default axis convention differs from the OpenSim and International Society of Biomechanics convention, and the converted model keeps the OpenSim convention: the axis remap is embedded in the body frames, joint values transfer unchanged, and only externally applied vectors such as ground-reaction forces require an explicit remap, which a pose-translation probe verified. The source model defines the patella through vastii insertion locations that move with knee angle rather than through a coordinate transform, and the converter realizes those moving insertions as coupled follower joints whose sampled range limits conflicted with the coupling constraints until the limits were removed; the rectus femoris route, which lacks the patella tracking, was rerouted over the vastii tracking point. Ground contact initially dragged on an invisible floor because the converted collision pairs ignore the per-geometry contact flags, and restricting contact to the six foot--ground pairs, with passive torques replacing the joint stops that the conversion dropped, removed the artifact. Comparing attachment paths with the OpenSim source checked the transferred geometry, and activating individual actuators checked the direction of the resulting joint torque. These checks concern conversion fidelity, not the ability of the controller to produce walking. The neural and mechanical models exchange signals at a \qty{2}{\milli\second} step.
